% !TEX root =./ECE444_Practice_main.tex
%
% L25 practice set (Tapering Lab). Lab-flavored: every question starts from a
% gain list read off the PHASER's Element Gains sliders or from a described
% sweep, and asks what the measurement will show. The recurring theme is the
% two-number discipline -- plotted peak drop vs taper efficiency.

\fullwidth{\textbf{Learning Objective 3.7: } Apply amplitude tapering to control
sidelobe level, and predict the resulting pattern trade-off.
\\[4pt]\normalfont\small Show your work. A \textbf{Documentation} line at the end
is required (write \textbf{None} if you did not collaborate).}

\begin{questions}

%%%%%%%%%%%%%%%%%%%%%%%%%%%%%%%%%%%%%%%%%%%%%%%%%%%%%%%%%%%%%%%%%%%%%%%%%%%%%
\question \textbf{Reading a taper off the sliders.} You find the eight Element
Gains sliders on the PHASER set to
\[ 30,\ 60,\ 90,\ 100,\ 100,\ 90,\ 60,\ 30 \ \ \text{percent.} \]
The array is the standard 8-element row, broadside, with the HB100 at boresight.
A uniform-taper sweep has already been frozen as the reference trace.

\begin{parts}

	\begin{minipage}[t]{0.58\linewidth}
		\part Compute the amplitude sum $\sum a_n$ and the peak drop this taper
		will show on the sweep plot relative to the frozen uniform trace.
		\begin{solution}[1.3in]
		The percentages are voltage amplitudes, so work in fractions of full scale:
		\eq{ \sum a_n &= 2\lp 0.30 + 0.60 + 0.90 + 1.00 \rp = 5.60 }
		The eight signals add in phase at boresight, so the summed voltage is
		proportional to $\sum a_n$ instead of $N = 8$:
		\eq{ \text{peak drop} &= 20\ \mylog{5.60/8} \\
		     &= 20\ \mylog{0.700} = -3.10\db }
		\end{solution}
	\end{minipage}\hfill%
	\begin{minipage}[t]{0.37\linewidth}
		\raggedleft \vspace{12pt}
		$\sum a_n$ = \ansbox[1.1in]{$5.60$}
		\\[10pt]
		peak drop = \ansbox[1.3in]{$-3.10\db$}
	\end{minipage}

	\begin{minipage}[t]{0.58\linewidth}
		\part Compute the taper efficiency $\eta_t$ for this distribution, as a
		ratio and in dB.
		\begin{solution}[1.3in]
		\eq{ \sum a_n^2 &= 2\lp 0.09 + 0.36 + 0.81 + 1.00 \rp = 4.52 \\
		     \eta_t &= \frac{\lp \sum a_n \rp^2}{N \sum a_n^2} = \frac{5.60^2}{8 \times 4.52} \\
		     &= \frac{31.36}{36.16} = 0.867 \quad \lp -0.62\db \rp }
		\end{solution}
	\end{minipage}\hfill%
	\begin{minipage}[t]{0.37\linewidth}
		\raggedleft \vspace{12pt}
		$\eta_t$ = \ansbox[1.1in]{$0.867$}
		\\[10pt]
		$\eta_t$ = \ansbox[1.1in]{$-0.62\db$}
	\end{minipage}

	\part Your two answers differ by 2.5 dB. State which one the Rectangular
	plot shows you, and which one belongs in a link budget as the antenna's loss.

		\begin{solution}[1.0in]
		The plot shows the $-3.10\db$ peak drop. That is the coherent
		receive-voltage loss, and it appears only because every trace is
		referenced to the same uniform-taper full scale. The link budget takes
		the $-0.62\db$ taper efficiency, because that is the directivity the
		array gave up. Charging a link budget the full 3.10 dB
		overstates the antenna's loss by 2.5 dB.
		\end{solution}

	\begin{minipage}[t]{0.58\linewidth}
		\part The uniform 8-element array has a broadside directivity of about
		8.9 dB. What is the directivity of the tapered array?
		\begin{solution}[1.0in]
		Directivity falls by the taper efficiency and by nothing else:
		\eq{ D &= 8.9\db - 0.62\db = 8.3\db }
		The peak drop plays no part in this calculation.
		\end{solution}
	\end{minipage}\hfill%
	\begin{minipage}[t]{0.37\linewidth}
		\raggedleft \vspace{12pt}
		$D$ = \ansbox[1.2in]{$8.3\dbi$}
	\end{minipage}

\end{parts}

\newpage
%%%%%%%%%%%%%%%%%%%%%%%%%%%%%%%%%%%%%%%%%%%%%%%%%%%%%%%%%%%%%%%%%%%%%%%%%%%%%
\question \textbf{Predicting a sweep.} A lab group presses the \textbf{Blackman}
preset, which sets the eight elements to
\[ 6,\ 27,\ 66,\ 100,\ 100,\ 66,\ 27,\ 6 \ \ \text{percent,} \]
and prepares to sweep. Their frozen uniform reference measured a
$13.1\mydeg$ half-power beamwidth with its first sidelobe at $-13$ dBc, and the
sweep's noise floor sits 23 dB below the uniform-taper peak.

\begin{parts}

	\begin{minipage}[t]{0.58\linewidth}
		\part Predict the peak drop the Blackman trace will show relative to the
		frozen uniform trace.
		\begin{solution}[1.2in]
		\eq{ \sum a_n &= 2\lp 0.06 + 0.27 + 0.66 + 1.00 \rp = 3.98 \\
		     \text{peak drop} &= 20\ \mylog{3.98/8} \\
		     &= 20\ \mylog{0.4975} = -6.1\db }
		\end{solution}
	\end{minipage}\hfill%
	\begin{minipage}[t]{0.37\linewidth}
		\raggedleft \vspace{12pt}
		peak drop = \ansbox[1.3in]{$-6.1\db$}
	\end{minipage}

	\begin{minipage}[t]{0.58\linewidth}
		\part The measured Blackman beamwidth is $23.1\mydeg$. By what factor has
		the beam broadened, and what physical change in the array produced it?
		\begin{solution}[1.2in]
		\eq{ \frac{23.1\mydeg}{13.1\mydeg} &= 1.76 }
		The taper leaves the end elements at 6\% of full amplitude, so the
		outer aperture contributes almost nothing. The array's effective
		aperture is smaller than its physical length, and beamwidth is set by
		aperture size.
		\end{solution}
	\end{minipage}\hfill%
	\begin{minipage}[t]{0.37\linewidth}
		\raggedleft \vspace{12pt}
		broadening = \ansbox[1.1in]{$1.76\times$}
	\end{minipage}

	\part Theory puts Blackman's first sidelobe roughly 33 dB below its own
	peak. What will the group see at that angle on the sweep, and what should
	they write in the sidelobe column of their table?

		\begin{solution}[1.1in]
		The Blackman peak sits 6.1 dB below the uniform peak, so a sidelobe 33 dB
		below it lands about 39 dB below full scale --- roughly 16 dB below the
		noise floor. Nothing is visible at that angle except the floor itself.
		The correct entry is \textbf{below the noise floor}. Reading a number off
		the grass would report the receiver's noise, not the antenna's sidelobe.
		\end{solution}

	\part A group member argues that Blackman is the better taper because its
	sidelobes are lower than Hann's. Using only what this measurement can show,
	evaluate that claim.

		\begin{solution}[1.2in]
		This measurement cannot support the claim. Hann's sidelobes are already
		below the noise floor, so both tapers produce the same observable
		result --- no visible sidelobes. Blackman costs $3.6\mydeg$ more
		beamwidth and 1.4 dB more peak drop than Hann, for an improvement the
		sweep cannot resolve. Against this receiver, Hann is the better choice;
		Blackman only earns its cost if the floor is lowered far enough to see
		Hann's sidelobes.
		\end{solution}

\end{parts}

\newpage
%%%%%%%%%%%%%%%%%%%%%%%%%%%%%%%%%%%%%%%%%%%%%%%%%%%%%%%%%%%%%%%%%%%%%%%%%%%%%
\question \textbf{A measurement that disagrees.} A student applies the
\textbf{Hann} preset ($12,\ 43,\ 77,\ 100,\ 100,\ 77,\ 43,\ 12$ percent),
predicts a $-4.7\db$ peak drop from the gain list, and measures $-6.8\db$
against the frozen uniform trace. The beam is also wider than the $19.5\mydeg$
the expectation table predicts.

\begin{parts}

	\begin{minipage}[t]{0.58\linewidth}
		\part Run the peak-drop formula backwards to find the amplitude sum the
		array achieved.
		\begin{solution}[1.2in]
		\eq{ \sum a_n &= N \times 10^{-6.8/20} = 8 \times 0.4571 = 3.66 }
		The commanded distribution sums to 4.64, so the array is short by
		$4.64 - 3.66 = 0.98$.
		\end{solution}
	\end{minipage}\hfill%
	\begin{minipage}[t]{0.37\linewidth}
		\raggedleft \vspace{12pt}
		$\sum a_n$ = \ansbox[1.1in]{$3.66$}
	\end{minipage}

	\part Identify the fault. Which element is misbehaving, and what amplitude
	is it contributing?

		\begin{solution}[1.0in]
		The shortfall of 0.98 is one full-amplitude element, so one of the two
		center elements (commanded to 100\%) is contributing nothing --- its
		gain is set to 0\%, or its channel has failed. No combination of the
		outer elements can account for a deficit that large: the two end
		elements together only carry 0.24.
		\end{solution}

	\begin{minipage}[t]{0.58\linewidth}
		\part Compute the taper efficiency of the array as it stands,
		with that element at zero.
		\begin{solution}[1.3in]
		The working elements are $0.12,\ 0.43,\ 0.77,\ 1.00,\ 0,\ 0.77,\ 0.43,\ 0.12$:
		\eq{ \sum a_n &= 3.64, \qquad \sum a_n^2 = 2.584 \\
		     \eta_t &= \frac{3.64^2}{8 \times 2.584} = \frac{13.25}{20.67} = 0.641 }
		which is $-1.9\db$, against $-1.2\db$ for the intended Hann taper.
		\end{solution}
	\end{minipage}\hfill%
	\begin{minipage}[t]{0.37\linewidth}
		\raggedleft \vspace{12pt}
		$\eta_t$ = \ansbox[1.1in]{$0.641$}
		\\[10pt]
		$\eta_t$ = \ansbox[1.1in]{$-1.9\db$}
	\end{minipage}

	\part Name one feature of the trace itself, other than the peak level, that
	confirms your diagnosis.

		\begin{solution}[1.0in]
		A dead center element breaks the symmetry of the distribution, so the
		measured pattern is no longer symmetric about boresight and the beam
		peak shifts slightly off $0\mydeg$. The gap in the middle of the
		aperture also raises the sidelobes back out of the noise floor, which a
		correctly applied Hann taper had buried.
		\end{solution}

\end{parts}

\newpage
%%%%%%%%%%%%%%%%%%%%%%%%%%%%%%%%%%%%%%%%%%%%%%%%%%%%%%%%%%%%%%%%%%%%%%%%%%%%%
\question \textbf{Designing to a specification.} You need a taper that keeps the
half-power beamwidth no wider than $17\mydeg$ while pushing the first sidelobe
below $-20$ dBc. With Enforce Symmetric Taper on, you set a linear ramp from
40\% at the ends to 91\% at the center:
\[ 40,\ 57,\ 74,\ 91,\ 91,\ 74,\ 57,\ 40 \ \ \text{percent.} \]
The measured sweep returns a $15.5\mydeg$ beamwidth with nothing visible above
the noise floor outside the main lobe.

\begin{parts}

	\begin{minipage}[t]{0.58\linewidth}
		\part Compute the peak drop and the taper efficiency of this design.
		\begin{solution}[1.4in]
		\eq{ \sum a_n &= 2\lp 0.40 + 0.57 + 0.74 + 0.91 \rp = 5.24 \\
		     \sum a_n^2 &= 2\lp 0.160 + 0.325 + 0.548 + 0.828 \rp = 3.721 \\
		     \text{peak drop} &= 20\ \mylog{5.24/8} = -3.7\db \\
		     \eta_t &= \frac{5.24^2}{8 \times 3.721} = 0.922 \quad \lp -0.35\db \rp }
		\end{solution}
	\end{minipage}\hfill%
	\begin{minipage}[t]{0.37\linewidth}
		\raggedleft \vspace{12pt}
		peak drop = \ansbox[1.2in]{$-3.7\db$}
		\\[10pt]
		$\eta_t$ = \ansbox[1.2in]{$0.922$}
	\end{minipage}

	\part The measurement shows nothing above the noise floor outside the main
	lobe. Explain what you can and cannot conclude about whether the $-20$ dBc
	part of the specification is met.

		\begin{solution}[1.2in]
		The floor sits 23 dB below the uniform peak, which is 19.3 dB below this
		trace's own peak. Nothing visible outside the main lobe therefore means
		the first sidelobe is at least 19.3 dB down. That demonstrates the
		specification is met to within the resolution of the measurement, but it
		does not produce a sidelobe number. To verify $-20$ dBc directly, the
		floor has to be lowered --- a stronger source, more averaging, or a
		quieter receiver.
		\end{solution}

	\begin{minipage}[t]{0.58\linewidth}
		\part Compare this design against the Hann preset ($-4.7\db$ peak drop,
		$-1.2\db$ taper efficiency, $19.5\mydeg$ beamwidth). State which you
		would fly on a radar that must resolve two targets $18\mydeg$ apart, and
		why.
		\begin{solution}[1.3in]
		The custom taper keeps $15.5\mydeg$ of beamwidth against Hann's
		$19.5\mydeg$, and gives up only 0.35 dB of directivity against Hann's
		1.2 dB. Two targets $18\mydeg$ apart fall inside a $19.5\mydeg$ beam and
		cannot be separated, so Hann fails the requirement outright. The custom
		taper resolves them and costs less directivity. Hann buys deeper
		sidelobes, but on this receiver that improvement is below the noise floor
		and cannot be measured.
		\end{solution}
	\end{minipage}\hfill%
	\begin{minipage}[t]{0.37\linewidth}
		\raggedleft \vspace{12pt}
		beamwidth margin = \ansbox[1.2in]{$2.5\mydeg$}
	\end{minipage}

	\part Enforce Symmetric Taper was on for this design. State what would
	change in the measured pattern if you set the same eight values with the
	switch off and mistyped one of them.

		\begin{solution}[1.0in]
		An asymmetric amplitude distribution produces an asymmetric pattern: the
		two sides of the trace no longer match, the sidelobes on one side rise
		above those on the other, and the beam peak moves off boresight. The
		symmetry switch makes that class of error impossible to enter by hand,
		which is why the lab keeps it on.
		\end{solution}

\end{parts}

\vspace{1.5em}
\noindent\textbf{Documentation:} \rule[-2pt]{0.6\linewidth}{0.4pt}

\end{questions}
